\subsection{CLOSURE OF AN EQUICONTINUOUS SET}

PROPOSITION 6. Let X be a topological (resp. uniform) space, let Y be a uniform space and let H be a subset of \(\mathcal{F}(X;Y)\) . Then H is equicontinuous at a point \(x_0 \in X\) (resp. uniformly equicontinuous) if and only if the closure \(\overline{\mathbf{H}}\) of H in \(\mathcal{F}_s(X;Y)\) is equicontinuous at \(x_0\) (resp. uniformly equicontinuous).

The condition is sufficient, trivially. To show that it is necessary, consider an entourage V of Y which is closed in \(Y \times Y\) ; by hypothesis, there is a neighbourhood U of \(x_{0}\) in X (resp. an entourage M of X) such that the relation \(x \in U\) (resp. \((x', x'') \in M\) ) implies \((h(x_{0}), h(x)) \in V\) {[}resp. \((h(x'), h(x'')) \in V\) {]} for all \(h \in H\) . Since V is closed, the mappings \(h \in \mathcal{F}(X; Y)\) which satisfy the relation \((h(x_{0}), h(x)) \in V\) for all \(x \in U\) {[}resp. the relation \((h(x'), h(x'')) \in V\) for all \((x', x'') \in M\) {]} form a closed subset of \(\mathcal{F}_{s}(X; Y)\) (§ 1, no. 2, Remark 6); since this closed subset contains H, it contains \(\overline{H}\) . Hence the result, since the closed entourages of Y form a fundamental system of entourages (Chapter II, § 1, no. 2, Proposition 2, Corollary 2).

